\scanpage{72}\mspage{71}
\deleted{\uline{\unsure{Cor. 4}}}

% vertical pen rule at the left of the definition
\noindent\llap{\smash{\rule[-1.7\baselineskip]{0.4pt}{1.6\baselineskip}}\enspace}\uline{\unsure{Definition|Définition}}\quad Une figure \deleted{\unsure{est} $F$} $\unsure{\Phi|\varphi}$ est dite
\uline{discrète} si toutes ses \deleted{\unsure{mult|multi}} strates sont
des lieux, i.e. $\widetilde\Phi\subset\unsure{\mathcal L|L}$.

\deleted{\uline{\unsure{Cor 4|Cor. 4}}}\quad \deleted{\unsure{Elle}}Une telle figure est connue
par la \unsure{connaissance} de la partie $\widetilde\Phi=A$
de $\mathcal L$. Celle-ci est \unsure{formée} de
lieux deux à deux disjoints, et \deleted{\unsure{pour}}
pour que une partie \uline{finie} de $\mathcal L$ ``soit''
une figure, il (f. et) s. qu'\unsure{elle|elles}
\deleted{\unsure{points|point} \unsure{soi|si|qui} \illegible{1} \unsure{soient}}
\inserted{soit \unsure{formée} de} lieux deux à deux
disjoints\unsure{..|.} \deleted{Plus \unsure{gén.|gén} , \unsure{sans|sa} condition de \unsure{finitude}.}
\inserted{\unsure{Notons|Notant|Noter} que la relation
d'ordre $\leq$ ou $\ll$ sur $\unsure{\widetilde\Phi|\overline{\widetilde\Phi}}=A$ est la relation}
discrète (c'est ainsi sur $\mathcal L$ !) donc
les parties fermées sont les parties
quelconques, donc les \unsure{s/figures|s.figures} \deleted{sont}
\unsure{correspondent} aux parties quelconques.

Plus \unsure{gén|gén.}\textsuperscript{t}, si $F$ est une figure,
on s'intéresse aux sous figures discrètes.
Elles sont décrites par \unsure{certaines} \deleted{$A\subset$}
parties $A$ de $\mathcal L$.

